\documentclass[11pt]{article}
\usepackage[utf8]{inputenc}
\usepackage[spanish,es-noindentfirst,es-noshorthands]{babel}
\usepackage[a4paper,margin=2cm]{geometry}
\usepackage[table]{xcolor}
\usepackage{array}
\usepackage{longtable}
\usepackage{enumitem}
\usepackage{fancyhdr}
\usepackage{amssymb}
\usepackage{hyperref}

\definecolor{vino}{HTML}{621F32}
\definecolor{oro}{HTML}{BC955C}
\definecolor{grisclaro}{HTML}{F1F5F9}
\definecolor{grisborde}{HTML}{CBD5E1}

\hypersetup{
  pdftitle={Caso de Prueba M6 - Valuacion Presupuestaria},
  pdfauthor={Eje Central - QA},
  colorlinks=true,
  linkcolor=vino,
  pdfborder={0 0 0}
}

\pagestyle{fancy}
\fancyhf{}
\renewcommand{\headrulewidth}{0.6pt}
\renewcommand{\footrulewidth}{0.3pt}
\fancyhead[L]{\small\color{vino}\textbf{CP-M06} --- Valuación Presupuestaria}
\fancyhead[R]{\small\color{grisborde} eje\_central\_front}
\fancyfoot[C]{\small\thepage}

\newcommand{\seccion}[1]{%
  \vspace{6pt}
  {\large\color{vino}\textbf{#1}}\\[2pt]
  {\color{oro}\rule{\linewidth}{1.2pt}}\\[6pt]
}

\newcommand{\subcaso}[2]{%
  \vspace{10pt}
  \noindent{\color{vino}\rule{\linewidth}{2.2pt}}\\[4pt]
  {\Large\color{vino}\textbf{#1}} \quad {\large\color{oro}\textbf{#2}}\\[4pt]
}

% Checkbox de PDF interactivo (AcroForm)
\newcommand{\chk}[1]{\CheckBox[name=#1,width=11pt,height=11pt,bordercolor=grisborde,backgroundcolor=1 1 1]{}}

% Campo de texto de una línea
\newcommand{\txtline}[3]{\TextField[name=#1,width=#2,height=13pt,bordercolor=grisborde,backgroundcolor=0.97 0.98 1]{#3}}

\begin{document}
\begin{Form}

% ---------- Portada / encabezado ----------
\begin{center}
  {\Huge\color{vino}\textbf{Caso de Prueba}}\\[4pt]
  {\Large\color{oro}\textbf{CP-M06}}\\[2pt]
  {\large Valuación Presupuestaria --- Simulador, Parámetros y Asuntos de Plazas}\\[2pt]
  {\small Módulo M6 --- 3 sub-casos: T1, T2, T3}
\end{center}

\vspace{6pt}
\noindent\small\textit{Instrucciones de uso: este documento no requiere conocimientos de
programación ni de sistemas. Ábrelo con un programa que permita rellenar formularios en PDF (Adobe
Acrobat Reader, Foxit, o el visor de Chrome/Edge) y marca cada casilla ($\Box$) conforme vayas
realizando cada paso con el mouse y el teclado, igual que lo haría cualquier persona que usa la
página. Todos los campos de texto y casillas son editables y se guardan al guardar el PDF. Si algún
paso no se puede completar, no te detengas: marca "Rechazado" o "Bloqueado" en el veredicto del
sub-caso y describe lo que pasó en "Observaciones".}\normalsize

\vspace{10pt}

\noindent
\begin{tabular}{|>{\columncolor{grisclaro}}p{3.2cm}|p{12.8cm}|}
\hline
\textbf{ID} & CP-M06 \\ \hline
\textbf{Módulo} & M6 --- Valuación Presupuestaria \\ \hline
\textbf{¿Dónde se hace la prueba?} & Iniciar sesión y, en el menú del lado izquierdo, entrar a
\emph{Valuación Presupuestaria}. Ahí aparecen 3 pestañas arriba: \emph{Simulador} (T1),
\emph{Parámetros} (T2) y \emph{Asuntos de Plazas} (T3). También se puede entrar escribiendo
directamente en la barra de direcciones del navegador: \path{/dashboard/valuacion_presupuestaria} \\ \hline
\textbf{Permisos necesarios} & Para usar \emph{Simulador} y \emph{Asuntos de Plazas} (T1 y T3) hace
falta el permiso \textbf{"Ver Valuación Presupuestaria"}. Para usar \emph{Parámetros} (T2) hace
falta el permiso \textbf{"Editar Parámetros de Valuación"}. Estos permisos los asigna un
administrador del sistema; si no se sabe si se cuenta con ellos, hay que preguntarle antes de
empezar. \\ \hline
\textbf{Importancia de esta prueba} & Alta --- este módulo calcula presupuesto oficial y permite
editar valores que afectan todos los cálculos del sistema \\ \hline
\textbf{Qué se va a probar} & Que los cálculos salgan correctos, que se pueda editar información
guardándose sola, que se puedan descargar reportes en PDF y Excel, y que cada pestaña solo la vea
quien tiene permiso para usarla \\ \hline
\textbf{Ejecutor} & \txtline{tester_nombre}{9cm}{} \\ \hline
\textbf{Fecha de ejecución} & \txtline{fecha_ejecucion}{5cm}{} \\ \hline
\textbf{Entorno} & \chk{env_local} Pruebas en mi equipo \quad \chk{env_staging} Pruebas (Staging) \quad \chk{env_prod} Producción (sistema real) \\ \hline
\end{tabular}

\seccion{Objetivo general}
Verificar que el módulo de Valuación Presupuestaria permite: (T1) calcular el impacto presupuestal
de un conjunto de plazas seleccionadas --- de forma manual o cargándolas automáticamente desde
plazas eventuales o permanentes ocupadas --- para un período de 1 a 12 meses, y descargar el
resultado en PDF y en Excel; (T2) editar directamente sobre la pantalla, guardándose de forma
automática, el catálogo de plazas (PECEN), los conceptos presupuestales y las constantes del
sistema que alimentan ese cálculo; y (T3) dar seguimiento a las solicitudes (asuntos) de ocupación
de plazas clasificadas como "Valuación Presupuestaria" y enviarlas al Simulador con su expediente
ya cargado. También se verifica que cada pestaña se muestre solo a quien tiene el permiso
correspondiente.

\seccion{Precondiciones generales}
\begingroup\small
\begin{itemize}[leftmargin=1.6em]
  \item[\chk{preg1}] Ya se inició sesión en el sistema.
  \item[\chk{preg2}] El usuario de prueba tiene AL MENOS UNO de los dos permisos:
    \textbf{"Ver Valuación Presupuestaria"} o \textbf{"Editar Parámetros de Valuación"}. Si no tiene
    ninguno de los dos, al entrar el sistema regresa automáticamente a la pantalla principal
    (Dashboard) sin mostrar ningún mensaje de "acceso denegado" --- esto es lo esperado, no es un
    error.
  \item[\chk{preg3}] Para ejecutar T1 y T3: el usuario tiene el permiso
    \textbf{"Ver Valuación Presupuestaria"}.
  \item[\chk{preg4}] Para ejecutar T2: el usuario tiene el permiso
    \textbf{"Editar Parámetros de Valuación"}.
  \item[\chk{preg5}] Existe en Control de Gestión al menos un oficio turnado con su asunto clasificado
    como "Valuación Presupuestaria", con expediente y al menos un documento adjunto (para T3 y para
    el flujo de integración T3 $\to$ T1).
  \item[\chk{preg6}] El catálogo de plazas (PECEN) tiene al menos un nivel activo con sueldo y
    compensación garantizada mayores a 0.
\end{itemize}
\endgroup

\noindent\textbf{Datos de prueba (llenar antes de ejecutar):}\\[4pt]
\begin{tabular}{|p{6.5cm}|p{9.5cm}|}
\hline
Usuario con permiso "Ver Valuación Presupuestaria" & \txtline{dato_user_view}{8.8cm}{} \\ \hline
Usuario con permiso "Editar Parámetros de Valuación" & \txtline{dato_user_edit}{8.8cm}{} \\ \hline
Usuario SIN ninguno de los dos permisos (para la prueba de acceso negado) & \txtline{dato_user_sin_permiso}{8.8cm}{} \\ \hline
Folio de asunto de prueba clasificado "Valuación Presupuestaria" & \txtline{dato_folio}{8.8cm}{} \\ \hline
Nivel de catálogo de prueba (para selección manual en el Simulador) & \txtline{dato_nivel}{8.8cm}{} \\ \hline
\end{tabular}

\seccion{Prueba de control de acceso (previa a los 3 sub-casos)}
\begingroup\small
\begin{longtable}{|c|p{6.7cm}|p{6.7cm}|c|}
\hline
\rowcolor{grisclaro}
\textbf{\#} & \textbf{Acción} & \textbf{Resultado esperado} & \textbf{OK} \\ \hline
\endhead

0.1 &
Iniciar sesión con el usuario que \textbf{no tiene ningún} permiso del módulo y escribir
directamente en la barra de direcciones del navegador: \path{/dashboard/valuacion_presupuestaria}. &
No se muestra ningún contenido del módulo ni mensaje de error; la página regresa automáticamente
a la pantalla principal (Dashboard). &
\chk{r01} \\ \hline

0.2 &
Iniciar sesión con el usuario que SOLO tiene el permiso "Ver Valuación Presupuestaria" (sin
"Editar Parámetros de Valuación") y entrar al módulo. &
La barra de pestañas muestra únicamente \emph{Simulador} y \emph{Asuntos de Plazas}; la pestaña
\emph{Parámetros} no aparece en absoluto (no es que esté apagada o bloqueada: simplemente no se
muestra). &
\chk{r02} \\ \hline

0.3 &
Iniciar sesión con el usuario que SOLO tiene el permiso "Editar Parámetros de Valuación" (sin
"Ver Valuación Presupuestaria") y entrar al módulo. &
La barra de pestañas muestra únicamente \emph{Parámetros}, y esta queda seleccionada
automáticamente (no existe \emph{Simulador} ni \emph{Asuntos de Plazas} para este usuario). &
\chk{r03} \\ \hline

\end{longtable}
\endgroup

% ======================================================================
\subcaso{CP-M06-T1}{Simulador de Valuación}
% ======================================================================

\textbf{Objetivo específico:} Verificar que el Simulador permite seleccionar plazas por nivel (de
forma manual o precargada desde "Eventuales Ocupadas"/"Permanentes Ocupadas"), calcular la
valuación presupuestaria para un período de 1 a 12 meses, mostrar el desglose correcto con sus
cuadros de ayuda con la fórmula usada, descargar el resultado en PDF y en Excel, y manejar el flujo
de "asunto vinculado" que llega
desde Asuntos de Plazas.

\vspace{4pt}\noindent\textbf{Precondición específica:} usar el usuario con el permiso "Ver
Valuación Presupuestaria". La pestaña que se muestra automáticamente al entrar al módulo es
\emph{Simulador}.

\begingroup\small
\begin{longtable}{|c|p{6.7cm}|p{6.7cm}|c|}
\hline
\rowcolor{grisclaro}
\textbf{\#} & \textbf{Acción} & \textbf{Resultado esperado} & \textbf{OK} \\ \hline
\endhead

1 &
Entrar a la pestaña \emph{Simulador}. &
Encabezado color vino con degradado y el texto "FUMP 2025 --- Sistema de Control de Plazas" /
"Valuación Presupuestaria"; debajo, un selector de período (12 botones, Ene a Dic); dos columnas:
"1 Selección de Niveles" (catálogo) y "2 Selección Activa" (vacía, con un ícono de capas y el texto
"Sin niveles seleccionados"); el botón "Calcular Valuación Presupuestaria" se ve apagado y no se
puede pulsar. &
\chk{t1r1} \\ \hline

2 &
Observar qué botón del selector de período queda resaltado en ámbar por defecto, sin hacer clic
en ninguno. &
Queda resaltado el botón cuya \textbf{posición numérica} (1 a 12, de izquierda a derecha) es igual
a (12 menos el número de mes actual contando enero = 0). Ese NÚMERO de posición --- no el nombre
del mes que muestra la etiqueta del botón --- es el período en meses que se usará por defecto al
calcular. &
\chk{t1r2} \\ \hline

3 &
Clic en cualquier otro botón del selector de período (ej. el último, "Dic"). &
El botón clicado queda resaltado en ámbar (fondo \#fbbf24, texto vino, ligeramente agrandado); el
que estaba resaltado antes pierde el resaltado; el período de evaluación pasa a ser la posición del
botón clicado (1 a 12). &
\chk{t1r3} \\ \hline

4 &
Escribir en el buscador "Filtrar por nivel o puesto..." (columna 1) un texto que coincida con el
nivel, código o denominación de al menos una fila. &
La lista de la columna 1 se filtra en vivo mostrando solo las filas cuyo nivel, código o
denominación contienen el texto (sin distinguir mayúsculas/minúsculas). Borrar el texto restaura
el catálogo completo. &
\chk{t1r4} \\ \hline

5 &
En una fila del catálogo, hacer clic 3 veces en la flecha hacia arriba ($\blacktriangle$) junto al
número de plazas. &
El número que aparece en el centro de esa fila sube de 0 a 3; al mismo tiempo aparece una tarjeta
con esa plaza en la columna 2 "Selección Activa" mostrando "3 plazas". &
\chk{t1r5} \\ \hline

6 &
En la misma fila, borrar el número escrito y escribir directamente "7". &
El campo queda en 7; la tarjeta correspondiente en "Selección Activa" actualiza su contador a
"7 plazas" (no se crea una tarjeta duplicada). &
\chk{t1r6} \\ \hline

7 &
En "Selección Activa", hacer clic en el ícono $\times$ de la plaza agregada en el paso anterior. &
Esa tarjeta desaparece de "Selección Activa"; el número de esa fila en el catálogo (columna 1)
regresa a vacío/0. &
\chk{t1r7} \\ \hline

8 &
Seleccionar 2 o más plazas distintas (repetir pasos 5--6 en otra fila) y luego hacer clic en el
botón "Limpiar" (arriba a la derecha de "Selección Activa"). &
Todas las plazas seleccionadas se eliminan de golpe; "Selección Activa" vuelve al estado vacío;
todos los números de plazas en el catálogo regresan a 0. El botón "Limpiar" solo se muestra cuando
hay al menos una plaza seleccionada. &
\chk{t1r8} \\ \hline

9 &
Clic en el botón "Eventuales Ocupadas". &
Mientras carga, el botón muestra un ícono girando y el texto "Cargando..."; tanto él como
"Permanentes Ocupadas" quedan apagados/no disponibles. Al terminar: "Eventuales Ocupadas" queda con
fondo vino sólido (activo); el catálogo y "Selección Activa" se llenan automáticamente con las
cantidades reales de plazas eventuales ocupadas que trae el sistema, sin que el usuario tenga que
escribir nada; cada fila con dato precargado muestra una etiqueta pequeña "N ocp." junto al nivel. &
\chk{t1r9} \\ \hline

10 &
Si el sistema encuentra niveles de nómina que no existen en el catálogo, observar el aviso bajo
los botones de carga. &
Aparece una franja roja con ícono de información y el texto "N nivel(es) de nómina sin
correspondencia en catálogo --- no incluidos en la carga.", donde N es la cantidad real que reporta
el sistema para el modo activo (Eventuales o Permanentes). &
\chk{t1r10} \\ \hline

11 &
Con "Eventuales Ocupadas" activo, hacer clic en "Permanentes Ocupadas". &
"Eventuales Ocupadas" se desactiva (pierde el fondo vino); "Permanentes Ocupadas" pasa a estar
activo con fondo azul (\#1a4a7a). Los datos de "Selección Activa" y del catálogo se
\textbf{reemplazan por completo} por los de plazas permanentes ocupadas (no se suman a los
anteriores). Ambos botones son mutuamente excluyentes: nunca quedan los dos activos a la vez. &
\chk{t1r11} \\ \hline

12 &
Con al menos una plaza con cantidad mayor a 0 seleccionada, hacer clic en "Calcular Valuación
Presupuestaria". &
El botón muestra un ícono girando y el texto "Procesando valuación..." y queda no disponible; al
recibir la respuesta, la página se desplaza sola y suavemente hasta la sección de resultados. &
\chk{t1r12} \\ \hline

13 &
Revisar la tabla "Desglose por Concepto". &
Encabezado vino degradado con etiqueta "FUMP 2025 $\cdot$ N Mes(es)" donde N coincide exactamente con el
período elegido en el paso 2 o 3; columnas Partida, Concepto, "Período Colectivo (Nm)",
"Regularizable (12m)", "Complemento Colectivo"; una fila por cada partida que trae el sistema;
fila final "Total Valuación Presupuestal" con fondo vino y montos en ámbar. &
\chk{t1r13} \\ \hline

14 &
Pasar el mouse (sin hacer clic) sobre cualquier celda de Partida, Concepto o de los 3 montos de
una fila de esa tabla. &
Aparece un cuadro de ayuda flotante oscuro con encabezado "Detalle de Cálculo" y, debajo, la
fórmula usada para ese concepto (ej. para la partida 12201: "$\Sigma$(Sueldo Base
$\times$ Plazas) $\times$ Meses"). El cuadro desaparece al quitar el mouse. &
\chk{t1r14} \\ \hline

15 &
Revisar la tabla "Desglose Analítico por Nivel". &
Encabezado con etiqueta "Base PECEN"; columnas Nivel, Plazas, Sueldo Base, "Sueldo Colectivo /
Período", "Comp. Garantizada", "Comp. Gar. Colectiva / Período", "Total Nivel"; una fila por cada
nivel seleccionado; fila "TOTAL" en fondo gris oscuro con montos en ámbar que suma todas las
filas. &
\chk{t1r15} \\ \hline

16 &
Clic en el botón "PDF" (arriba a la derecha de la sección de resultados). &
Se descarga un archivo \path{Valuacion_Presupuestaria_<timestamp>.pdf} con encabezado vino
"REPORTE DE VALUACIÓN PRESUPUESTARIA", fecha de generación, período evaluado y --- solo si se usó
"Eventuales Ocupadas" o "Permanentes Ocupadas" --- una etiqueta con el modo usado; contiene ambas
tablas (Desglose por Concepto y Desglose Analítico por Nivel) con sus totales, más una fila
adicional "QUINCENA" en la segunda tabla. &
\chk{t1r16} \\ \hline

17 &
Clic en el botón "Excel". &
Se descarga \path{Valuacion_Presupuestaria_<timestamp>.xlsx} con 2 hojas: "Resumen por
Concepto" y "Desglose Analítico"; ambas con banda de título vino, encabezados en vino, montos con
formato de moneda, fila TOTAL y (solo en la hoja "Desglose Analítico") fila QUINCENA adicional. &
\chk{t1r17} \\ \hline

18 &
Ir a la pestaña \emph{Asuntos de Plazas} (sub-caso T3) y hacer clic en el botón "Ir al simulador"
(ícono de calculadora) de cualquier fila. &
La vista cambia automáticamente a la pestaña \emph{Simulador}; el diseño cambia de 2 columnas a un
panel dividido: panel izquierdo fijo (48\% del ancho) con encabezado "Oficio Vinculado" + folio
del asunto + botón "Cerrar Asunto"; visor de PDF del primer documento del expediente (o el mensaje
"No se pudo cargar el documento" si no existe ninguno); bloque de metadatos (Número de Oficio,
Remitente, Descripción) y lista de "Archivos del Expediente" / "Respuestas" si existen. Panel
derecho (52\%) con el Simulador completo, igual al de los pasos 1--17. &
\chk{t1r18} \\ \hline

19 &
Si el expediente tiene más de un documento/respuesta listados, hacer clic en uno distinto al que
se está mostrando. &
Mientras carga se ve una pantalla con el aviso "Generando Vista Previa..." (o un ícono girando);
al terminar, el visor de PDF cambia al documento seleccionado y el botón de ese documento queda
resaltado (vino para documentos, verde para respuestas). &
\chk{t1r19} \\ \hline

20 &
Clic en "Cerrar Asunto". &
El diseño regresa al de 2 columnas normal (sin panel de expediente); la selección de plazas y el
resultado de cálculo obtenidos en los pasos previos \textbf{no se pierden} (persisten en pantalla). &
\chk{t1r20} \\ \hline

\end{longtable}
\endgroup

\noindent\textbf{Criterios de aceptación T1:}
\begin{itemize}[leftmargin=1.6em]
  \item[\chk{t1ac1}] "Calcular Valuación Presupuestaria" permanece deshabilitado mientras no haya
    ninguna plaza con cantidad mayor a 0.
  \item[\chk{t1ac2}] "Eventuales Ocupadas" y "Permanentes Ocupadas" son mutuamente excluyentes;
    activar uno reemplaza por completo los datos cargados por el otro.
  \item[\chk{t1ac3}] El período (1--12) usado en el cálculo coincide exactamente con el botón
    resaltado en el selector al momento de pulsar "Calcular".
  \item[\chk{t1ac4}] Los totales mostrados en pantalla ("Total Valuación Presupuestal" y "TOTAL"
    por nivel) coinciden exactamente con los totales exportados a PDF y a Excel.
  \item[\chk{t1ac5}] El flujo de "Asunto vinculado" no resetea una selección de plazas hecha
    previamente en el Simulador.
  \item[\chk{t1ac6}] Durante todo el flujo no aparecen mensajes de error inesperados, pantallas en
    blanco, ni el sistema se queda congelado o sin responder.
\end{itemize}

\vspace{6pt}\noindent
\textbf{Veredicto T1:}\quad
\chk{t1_aprobado} \textbf{Aprobado} \qquad
\chk{t1_rechazado} \textbf{Rechazado} \qquad
\chk{t1_bloqueado} \textbf{Bloqueado}\\[4pt]
\noindent\textbf{Observaciones T1:}\\[4pt]
\TextField[name=obs_t1, width=16cm, height=1.6cm, multiline=true, bordercolor=grisborde, backgroundcolor=0.97 0.98 1]{}

% ======================================================================
\subcaso{CP-M06-T2}{Parámetros de Valuación}
% ======================================================================

\textbf{Objetivo específico:} Verificar que la pestaña \emph{Parámetros} permite editar
directamente sobre la pantalla, guardándose de forma automática y real, las 3 tablas que alimentan
el cálculo de valuación (Catálogo de Plazas, Conceptos Presupuestales, Constantes), y que los
filtros/búsqueda funcionan y se reinician correctamente al cambiar de sub-pestaña. Este sub-caso
solo aplica a usuarios con el permiso "Editar Parámetros de Valuación".

\vspace{4pt}\noindent\textbf{Precondición específica:} usar el usuario con el permiso "Editar
Parámetros de Valuación". Entrar a la pestaña \emph{Parámetros}.

\begingroup\small
\begin{longtable}{|c|p{6.7cm}|p{6.7cm}|c|}
\hline
\rowcolor{grisclaro}
\textbf{\#} & \textbf{Acción} & \textbf{Resultado esperado} & \textbf{OK} \\ \hline
\endhead

1 &
Clic en la pestaña \emph{Parámetros}. &
Se muestran 3 sub-pestañas: "Catálogo Plazas", "Conceptos Pres." y "Constantes" (seleccionada por
defecto: "Catálogo Plazas"). El encabezado de contenido muestra el título "Catálogo PECEN" (nota: el
título de la página NO es idéntico al texto del botón de la sub-pestaña) junto con la leyenda
"Sistema FUMP $\cdot$
Clic en valor para editar $\cdot$ Enter para guardar"; a la derecha: buscador "Filtrar...", etiqueta verde
"$\checkmark$ Auto guardado", botón "Reiniciar filtros" (deshabilitado, sin filtros activos) y botón
"Recargar". &
\chk{t2r1} \\ \hline

2 &
Sin tocar nada más, revisar el pie fijo de la tabla. &
Franja ámbar con ícono de alerta y el texto exacto: "Precaución: Las modificaciones afectarán
directamente todos los cálculos de valuación en tiempo real." Esta franja está presente en las 3
sub-pestañas. &
\chk{t2r2} \\ \hline

3 &
En "Catálogo Plazas", clic en el valor de la columna "Sueldo" de cualquier fila. &
La celda entra en modo edición: aparece un campo de texto con borde inferior vino y fondo ámbar
claro, ya con el valor actual escrito. &
\chk{t2r3} \\ \hline

4 &
Escribir un número nuevo distinto y presionar Enter. &
La celda sale de edición mostrando el nuevo valor ya formateado (2 a 4 decimales, formato es-MX);
aparece brevemente un pequeño ícono girando junto al valor mientras guarda; al recargar la página
(tecla F5 o el botón de recargar del navegador) el nuevo valor sigue ahí (esto confirma que sí se
guardó). &
\chk{t2r4} \\ \hline

5 &
Clic en otra celda editable (ej. "Denominación"), escribir un valor distinto y presionar Escape
antes de hacer clic fuera. &
La edición se cancela y la celda regresa a mostrar su valor original; no se guarda nada (al
recargar la página se confirma que el valor no cambió). &
\chk{t2r5} \\ \hline

6 &
Clic en una celda editable, cambiar el valor y hacer clic fuera de la celda sin presionar Enter. &
El cambio se guarda igual que con Enter --- el guardado automático también funciona al hacer clic
fuera del campo, no solo al presionar Enter. &
\chk{t2r6} \\ \hline

7 &
Editar una celda numérica y forzar que el guardado falle (ej. cortar la conexión a internet antes
de salir del campo, o escribir un valor que el sistema rechace). &
Aparece una alerta nativa del navegador con el texto "Error al guardar el cambio"; al cerrar la
alerta, la celda regresa a mostrar su valor original (no se queda con el valor inválido). &
\chk{t2r7} \\ \hline

8 &
Escribir en el buscador global "Filtrar..." un texto que coincida con el nivel o denominación de
al menos una fila. &
La tabla se reduce en vivo a las filas coincidentes; el botón "Reiniciar filtros" pasa de
deshabilitado a habilitado. &
\chk{t2r8} \\ \hline

9 &
Clic en el ícono de embudo del encabezado de una columna (ej. "Zona"). &
Se abre un menú desplegable con una lista de casillas con los valores únicos de esa columna
(igual que en Excel); al deseleccionar
algunos valores y aplicar, la tabla se reduce a las filas cuyo valor de esa columna quedó marcado;
el ícono de embudo de esa columna queda resaltado en vino con un punto indicador. &
\chk{t2r9} \\ \hline

10 &
Clic en "Reiniciar filtros" (con buscador y/o filtro de columna activos). &
Buscador global, filtros de columna y filtros de texto se limpian todos a la vez; la tabla regresa
a mostrar el catálogo completo; el botón vuelve a quedar deshabilitado. &
\chk{t2r10} \\ \hline

11 &
Clic en "Recargar". &
Se descartan los cambios que aún no se reflejaban en pantalla y se vuelve a traer el catálogo, los
conceptos y las constantes desde el sistema (mientras dura, se ve una pantalla de carga que cubre
toda la vista). &
\chk{t2r11} \\ \hline

12 &
Cambiar a la sub-pestaña "Conceptos Pres.". &
El encabezado cambia a "Conceptos Presupuestales"; tabla con columnas Partida, Descripción,
Sección, Orden; los filtros de columna y el texto del buscador que hubiera quedado activo en
"Catálogo Plazas" se resetean (no se conservan al cambiar de sub-pestaña); la columna "Partida" no es
editable (texto mono fijo en vino), el resto sí. &
\chk{t2r12} \\ \hline

13 &
Cambiar a la sub-pestaña "Constantes". &
El encabezado cambia a "Constantes del Sistema"; tabla con columnas Clave, "Descripción
Operativa", Valor; "Clave" no es editable; "Descripción Operativa" y "Valor" sí, con el mismo
patrón: clic para editar, Enter o clic fuera para guardar, Escape para cancelar. &
\chk{t2r13} \\ \hline

\end{longtable}
\endgroup

\noindent\textbf{Criterios de aceptación T2:}
\begin{itemize}[leftmargin=1.6em]
  \item[\chk{t2ac1}] Ningún cambio se guarda hasta hacer clic fuera del campo o presionar Enter;
    Escape siempre cancela sin guardar nada.
  \item[\chk{t2ac2}] Cambiar de sub-pestaña (Catálogo / Conceptos / Constantes) limpia el buscador y
    los filtros de columna de la sub-pestaña anterior.
  \item[\chk{t2ac3}] La pestaña "Parámetros" no se muestra en absoluto (no solo aparece apagada) para
    un usuario sin el permiso "Editar Parámetros de Valuación".
  \item[\chk{t2ac4}] El aviso de precaución sobre impacto en cálculos en tiempo real está presente
    en las 3 sub-pestañas.
  \item[\chk{t2ac5}] Durante todo el flujo no aparecen mensajes de error inesperados, pantallas en
    blanco, ni el sistema se queda congelado o sin responder.
\end{itemize}

\vspace{6pt}\noindent
\textbf{Veredicto T2:}\quad
\chk{t2_aprobado} \textbf{Aprobado} \qquad
\chk{t2_rechazado} \textbf{Rechazado} \qquad
\chk{t2_bloqueado} \textbf{Bloqueado}\\[4pt]
\noindent\textbf{Observaciones T2:}\\[4pt]
\TextField[name=obs_t2, width=16cm, height=1.6cm, multiline=true, bordercolor=grisborde, backgroundcolor=0.97 0.98 1]{}

% ======================================================================
\subcaso{CP-M06-T3}{Asuntos de Plazas}
% ======================================================================

\textbf{Objetivo específico:} Verificar que la pestaña \emph{Asuntos de Plazas} lista únicamente las
solicitudes de ocupación de plazas clasificadas como "Valuación Presupuestaria" que además tienen
un oficio turnado coincidente en Control de Gestión, que el buscador y las acciones "Ver detalles"
e "Ir al simulador" funcionan correctamente, y que reclasificar un asunto lo remueve de esta vista.

\vspace{4pt}\noindent\textbf{Precondición específica:} usar el usuario con el permiso "Ver
Valuación Presupuestaria". Entrar a la pestaña \emph{Asuntos de Plazas}.

\begingroup\small
\begin{longtable}{|c|p{6.7cm}|p{6.7cm}|c|}
\hline
\rowcolor{grisclaro}
\textbf{\#} & \textbf{Acción} & \textbf{Resultado esperado} & \textbf{OK} \\ \hline
\endhead

1 &
Clic en la pestaña \emph{Asuntos de Plazas}. &
Encabezado color vino con el texto "Asuntos de Solicitud de Ocupación de Plazas" y subtítulo
"Dirección de Organización - Valuación Presupuestaria"; contador "Total de Solicitudes" con
animación de conteo ascendente. Mientras carga: el contador muestra un bloque gris parpadeante y la
tabla muestra 4 filas grises animadas a modo de vista previa (todavía no es la tabla real). &
\chk{t3r1} \\ \hline

2 &
Esperar a que termine de cargar. &
La tabla muestra solo asuntos que tienen un oficio turnado coincidente en Control de Gestión (los
que no lo tienen quedan excluidos sin generar error visible ni fila vacía); columnas: "Folio /
Oficio", "Remitente / Dependencia", "Descripción Asunto", "Valuación", "Acciones"; pie de tabla con
el texto "Mostrando N solicitudes". &
\chk{t3r2} \\ \hline

3 &
Revisar la columna "Valuación" de varias filas con distinto estatus. &
Etiqueta de color según el estatus: verde "Procedente", rojo "Improcedente", ámbar "Pendiente"
(este último es el valor que se muestra cuando no hay estatus o no coincide con los dos
anteriores). Si el registro tiene una liga de resolución cargada, debajo de la etiqueta aparece el
enlace "Ver Resolución" (ícono de clip) que abre esa dirección en una pestaña nueva del navegador. &
\chk{t3r3} \\ \hline

4 &
Escribir en el buscador "Buscar por folio, oficio o remitente..." un texto que coincida con folio,
número de oficio, remitente o descripción de una fila conocida. &
La tabla se filtra en vivo (no consulta al sistema, es un filtrado que ocurre en la misma pantalla); el pie actualiza
"Mostrando N solicitudes" al nuevo total filtrado. &
\chk{t3r4} \\ \hline

5 &
Escribir un texto que no coincida con ninguna fila existente. &
La tabla muestra un estado vacío centrado: ícono de embudo, texto "Sin resultados de solicitudes" y
--- solo porque hay texto en el buscador --- un botón "Limpiar filtro". &
\chk{t3r5} \\ \hline

6 &
Clic en "Limpiar filtro" (o en el ícono $\times$ dentro del buscador). &
El buscador se vacía y la tabla vuelve a mostrar todas las solicitudes clasificadas. &
\chk{t3r6} \\ \hline

7 &
Clic en el botón de acción "Ver detalles" (ícono de ojo) de una fila. &
Se abre una ventana emergente con el detalle (la misma que se usa en Oficios Turnados a DO).
Mientras carga el expediente o documento se ve el aviso "Generando Vista Previa...". La ventana
muestra el folio y el número de oficio en el título, una etiqueta con el estatus del oficio
turnado, un selector "Clasificar:" con el tipo de asunto actual, botones "Anterior" / "Siguiente"
para navegar, y un botón para cerrar (X). &
\chk{t3r7} \\ \hline

8 &
Con la ventana abierta y más de un asunto en la lista ya filtrada, hacer clic en el botón
"Siguiente". &
La ventana muestra los datos del siguiente asunto de la LISTA YA FILTRADA (si había texto en el
buscador al abrir la ventana, avanza dentro de ese subconjunto, no de la lista completa). En el
último elemento de esa lista, "Siguiente" queda apagado; en el primero, "Anterior" queda apagado. &
\chk{t3r8} \\ \hline

9 &
Cerrar la ventana con el botón (X). &
La ventana se cierra; la tabla de fondo conserva el mismo texto de búsqueda que tenía antes de
abrirla. &
\chk{t3r9} \\ \hline

10 &
Clic en el botón de acción "Ir al simulador" (ícono de calculadora) de una fila. &
La vista cambia automáticamente a la pestaña \emph{Simulador} con el expediente de esa fila
precargado en el panel izquierdo dividido (ver paso 18 de T1). &
\chk{t3r10} \\ \hline

11 &
Regresar a la pestaña \emph{Asuntos de Plazas}, abrir "Ver detalles" de un asunto de prueba y, en
el selector "Clasificar:" de la ventana, cambiar la clasificación a un tipo DISTINTO de "Valuación
Presupuestaria", confirmar y cerrar la ventana. &
Mientras guarda se ve un pequeño ícono girando junto al selector. Al cerrar la ventana, la lista de
Asuntos de Plazas se recarga automáticamente y ese asunto YA NO aparece en la tabla (dejó de estar
clasificado como Valuación Presupuestaria). &
\chk{t3r11} \\ \hline

\end{longtable}
\endgroup

\noindent\textbf{Criterios de aceptación T3:}
\begin{itemize}[leftmargin=1.6em]
  \item[\chk{t3ac1}] Ningún asunto sin oficio turnado coincidente aparece en la tabla ni genera
    error visible en pantalla.
  \item[\chk{t3ac2}] El contador "Total de Solicitudes" (encabezado) y el pie "Mostrando N
    solicitudes" reflejan, respectivamente, el total sin filtrar y el total ya filtrado por el
    buscador.
  \item[\chk{t3ac3}] Reclasificar un asunto fuera de "Valuación Presupuestaria" lo remueve de esta
    vista tras el refresco automático.
  \item[\chk{t3ac4}] La navegación Anterior/Siguiente de la ventana de detalle respeta la lista ya
    filtrada por el buscador, no la lista completa sin filtrar.
  \item[\chk{t3ac5}] Durante todo el flujo no aparecen mensajes de error inesperados, pantallas en
    blanco, ni el sistema se queda congelado o sin responder.
\end{itemize}

\vspace{6pt}\noindent
\textbf{Veredicto T3:}\quad
\chk{t3_aprobado} \textbf{Aprobado} \qquad
\chk{t3_rechazado} \textbf{Rechazado} \qquad
\chk{t3_bloqueado} \textbf{Bloqueado}\\[4pt]
\noindent\textbf{Observaciones T3:}\\[4pt]
\TextField[name=obs_t3, width=16cm, height=1.6cm, multiline=true, bordercolor=grisborde, backgroundcolor=0.97 0.98 1]{}

% ======================================================================
\seccion{Resultado global del módulo M6}
% ======================================================================

\noindent
\begin{tabular}{|p{5.5cm}|p{4.7cm}|p{4.7cm}|}
\hline
\rowcolor{grisclaro}
\textbf{Sub-caso} & \textbf{Veredicto} & \textbf{Referencia} \\ \hline
CP-M06-T1 --- Simulador & Ver casilla marcada en T1 & Sección "Veredicto T1" \\ \hline
CP-M06-T2 --- Parámetros & Ver casilla marcada en T2 & Sección "Veredicto T2" \\ \hline
CP-M06-T3 --- Asuntos de Plazas & Ver casilla marcada en T3 & Sección "Veredicto T3" \\ \hline
\end{tabular}

\vspace{10pt}\noindent
\textbf{Veredicto general del módulo:}\quad
\chk{veredicto_aprobado} \textbf{Aprobado} \qquad
\chk{veredicto_rechazado} \textbf{Rechazado} \qquad
\chk{veredicto_bloqueado} \textbf{Bloqueado}

\vspace{8pt}
\noindent\textbf{Observaciones generales / defectos encontrados:}\\[4pt]
\TextField[name=observaciones, width=16cm, height=3cm, multiline=true, bordercolor=grisborde, backgroundcolor=0.97 0.98 1]{}

\vspace{10pt}
\noindent\textbf{Evidencia adjunta (rutas de captura/video):}\\[4pt]
\txtline{evidencia}{16cm}{}

\end{Form}
\end{document}
